\chapter{Conclusion}
\section{Summary of the thesis}
\label{sec:summary_thesis}

This thesis has investigated whether, and to what extent, the framework introduced in \emph{Universal Quantum Control through Deep Reinforcement Learning} can be regarded as a physically meaningful and empirically credible strategy for robust analog quantum control on superconducting qubits \cite{Niu2019}. More specifically, the work was structured around a double objective. The first was to reconstruct, as faithfully as possible, the main physical and methodological ingredients of the original framework, including the multilevel gmon model, the leakage-aware UFO objective, and the continuous-action reinforcement-learning formulation based on TRPO. The second was to assess this framework critically through an independent empirical study, with particular attention to its nominal performance, its robustness under stochastic perturbations, and its runtime behavior across a physically relevant family of target gates.

From the theoretical point of view, the thesis first developed the broader control-theoretic and physical background needed to situate the problem properly. In particular, it emphasized that realistic quantum gate design on near-term superconducting hardware is intrinsically a multi-objective problem, in which gate fidelity, runtime, leakage suppression, and robustness cannot be treated independently. This perspective motivated the discussion of weakly anharmonic multilevel qubits, computational and leakage subspaces, the speed--leakage tradeoff, and the general role of model-based optimal control and reinforcement learning in pulse-level quantum control. On this basis, the thesis then reconstructed the paper-specific framework in a more explicit and organized form, clarifying the role of the effective gmon Hamiltonian, the time-dependent Schrieffer--Wolff treatment of leakage, the structure of the UFO cost function, and the logic by which the control problem is cast as a sequential decision process.

From the implementation point of view, the thesis provided a coherent simulation pipeline in which the original framework could be studied quantitatively. The controlled dynamics were simulated in a truncated multilevel Hilbert space, the target gates were chosen from the family \(N(\alpha,\alpha,\gamma)\), and the control objective was evaluated through the same general balance of fidelity, leakage, boundary regularity, and runtime that underlies the original UFO proposal. This made it possible to compare different optimization strategies under matched physical assumptions and to interpret the results within a common methodological setting, rather than as isolated numerical exercises.

The empirical part of the thesis then led to a differentiated but overall coherent picture. On the representative single-target benchmark, the direct Adam baseline proved stronger than the nominal TRPO controller in the deterministic regime, showing that reinforcement learning does not automatically dominate a strong direct pulse optimizer on a fixed nominal problem. However, once the analysis was extended to stochastic control perturbations, the noise-optimized TRPO controller emerged as the strongest overall method, achieving both the highest average fidelity and the lowest fidelity variance under noisy evaluation. This result supports one of the main scientific messages of the original paper: robustness against stochastic control errors can be learned more effectively when the relevant perturbation model is already embedded in the training environment. In parallel, the runtime sweep over the family \(N(\alpha,\alpha,\pi/2)\) showed that analog control can remain substantially faster than standard gate synthesis over a broad region of target space, with the strongest gains appearing where the target gate is most naturally aligned with the native interaction structure of the hardware.

Taken together, these results suggest that the main significance of the validated framework lies not in offering a universally dominant nominal optimizer, but in providing a unified and physically well-motivated strategy for robust analog quantum control. The present thesis therefore supports the view that the real strength of the UFO framework is its ability to integrate, within a single control architecture, several competing requirements that are central to realistic superconducting-qubit operation. In this sense, the work contributes both a structured reconstruction of the original framework and a more critical empirical assessment of its range of validity, its practical advantages, and its limitations.

\textbf{TODO:} Inserire qui, se opportuno, una frase finale che richiami in modo sintetico i risultati definitivi dello studio sulle architetture delle policy/value network.

\textbf{TODO:} Valutare se aggiungere un breve richiamo conclusivo ai risultati definitivi della sensitivity analysis dei pesi della UFO cost function, solo se tali risultati verranno ritenuti rilevanti per il messaggio complessivo della tesi.


\section{Limitations and future work}
\label{sec:limitations_future_work}

Despite the encouraging results obtained in this thesis, the present study also has a number of clear limitations that should be stated explicitly. First, the entire analysis is simulation-based. Although the physical model is designed to remain close to a realistic superconducting-qubit setting, the framework has not been validated on actual hardware. As a consequence, the results should be interpreted as evidence about the behavior of the method within a physically motivated numerical environment, rather than as a direct demonstration of experimental performance. In particular, real devices may introduce additional distortions, calibration drifts, transfer-function nonlinearities, and measurement constraints that are only partially captured by the present simulation pipeline.

A second limitation concerns the modeling assumptions used throughout the work. The dynamics are simulated in a truncated multilevel Hilbert space with three local levels per mode, under the effective rotating-wave description adopted in the original framework. While this is sufficient to capture the dominant leakage channels and to preserve the main physical structure of the problem, it remains an approximation. Higher excited states, more detailed control-transfer effects, and more complete open-system dynamics may all become relevant in experimentally realistic regimes. Similarly, the stochastic robustness study has been carried out primarily under a Gaussian control-noise model, with fluctuations injected into the control amplitudes and into the anharmonicity-related parameter. This is physically meaningful and consistent with the original paper, but it does not exhaust the broader space of possible imperfections affecting superconducting hardware.

A third limitation concerns the optimization comparison itself. Although the thesis includes a meaningful non-RL baseline through direct Adam-based pulse optimization, and therefore does not compare the RL controller against a weak reference, the space of alternative control methods is much broader. The present study does not provide a systematic matched comparison against the full family of established quantum optimal-control approaches, such as more refined GRAPE variants, GOAT-like parametrizations, Krotov-type methods, or hybrid open-loop / closed-loop calibration strategies. For this reason, the conclusions of the thesis should be understood as a comparison between a specific RL-based framework and a specific strong direct-optimization baseline, rather than as a definitive ranking of all available control methodologies.

A further limitation is that some parts of the empirical assessment remain sensitive to implementation-level choices. As discussed in the parameter-study section, quantities such as the training horizon, stopping criterion, policy/value architecture, and the scalarization induced by the UFO weights can affect the quantitative behavior of the learned controller. This does not invalidate the main qualitative conclusions of the thesis, but it does imply that some numerical details of the results should be interpreted with caution. In particular, the strongest conclusions of the work concern the qualitative role of stochastic training and the existence of a runtime advantage in hardware-aligned regions of target space, rather than the exact numerical value of every observed improvement.

These limitations suggest several natural directions for future work. A first priority would be to extend the robustness study beyond the Gaussian perturbation model used here. In particular, it would be scientifically valuable to test whether the robustness learned under stochastic amplitude noise generalizes to quasi-static drift, temporally correlated noise, parameter mismatch, and simple open-system decoherence channels. Such an extension would make it possible to distinguish more clearly between robustness that is specific to the training distribution and robustness that reflects a broader structural regularity of the learned control strategy.

A second important direction concerns the comparison with more standard quantum optimal-control techniques. The present results already show that direct gradient-based optimization can remain very strong in the nominal single-target setting. A more systematic future study could therefore compare RL-based control, GRAPE-like methods, and low-parameter analytic approaches under fully matched assumptions, keeping the same Hamiltonian model, truncation level, control bounds, runtime penalty, and robustness metrics. This would help clarify more sharply which part of the observed performance comes from the physical objective itself, which part comes from the optimization strategy, and which part comes from the interaction between the two.

A third natural direction is to deepen the analysis of the runtime landscape observed on the family \(N(\alpha,\alpha,\pi/2)\). The present sweep suggests that the analog-control advantage is strongly structured, with broad favorable basins, harder edge regions, and isolated points of especially strong alignment with the native hardware dynamics. This raises interesting open questions about the geometry of the control landscape, the origin of the apparent singular features, and the role of curriculum transfer in stabilizing short-time solutions across nearby targets. A more detailed theoretical and numerical study of these aspects could improve the interpretability of the framework and clarify whether the observed runtime structure is specific to this target family or reflects a more general property of analog quantum control on weakly anharmonic superconducting platforms.

More broadly, an especially promising long-term direction would be to move from purely simulation-based open-loop optimization toward hybrid workflows that combine model-based learning with hardware-level refinement. One of the most attractive possibilities is to use the present framework to generate physically informed robust initial pulses, and then refine them through experimental calibration or partial closed-loop updates on real devices. In this way, one could attempt to combine the interpretability and structure of the UFO objective with the realism and adaptability of hardware-in-the-loop optimization.

In summary, the present thesis supports the framework of robust analog quantum control through reinforcement learning as a physically meaningful and promising strategy, but it also shows that its strengths are most convincing when interpreted with appropriate precision. Its main value lies in the robust and hardware-aware regime, not in a universal nominal dominance over all alternatives. Future work should therefore focus on broadening the robustness analysis, sharpening the comparison with established control methods, and testing how well the framework survives the transition from a carefully structured simulation environment to increasingly realistic experimental conditions.

\textbf{TODO:} Inserire qui una breve limitazione finale legata ai risultati definitivi dello studio sulle architetture delle policy/value network, una volta completata la versione conclusiva di quella analisi.

\textbf{TODO:} Valutare se aggiungere qui una frase specifica sui possibili sviluppi futuri relativi alla calibrazione o reinterpretazione dei pesi della UFO cost function, dopo il consolidamento definitivo della sensitivity analysis.